%% Согласно ГОСТ Р 7.0.11-2011:
%% 5.3.3 В заключении диссертации излагают итоги выполненного исследования, рекомендации, перспективы дальнейшей разработки темы.
%% 9.2.3 В заключении автореферата диссертации излагают итоги данного исследования, рекомендации и перспективы дальнейшей разработки темы.
\begin{enumerate}
  \item проведён анализ методов безопасного обучения с подкреплением и выявлены ограничения одноуровневых стратегий безопасного обучения с подкреплением, безопасного обучения с подкреплением на основе модели мира и существующих методов безопасного иерархического обучения с подкреплением;
  \item исследована роль модели мира в оценке награды, штрафов и долгосрочного риска при безопасном планировании;
  \item разработан метод SG-Safe для безопасного исследования в задачах с длинным горизонтом за счёт генерации промежуточных подцелей;
  \item разработан метод ITES, объединяющий иерархическую декомпозицию, модель мира и модель штрафов для контроля безопасности на уровне подцелей и действий;
  \item разработана мультимодальная постановка безопасного обучения с подкреплением с естественно-языковыми целями и ограничениями и проведена экспериментальная оценка современных агентов в этой постановке.
\end{enumerate}
